% GENERATED FILE. Edit canonical lessons, then run npm run generate:books.
% canonical-source-hash: fnv1a64:f76f941cd2c0e492
% canonical-lessons: MR-C303-ekekali, MR-C303-katha, MR-C303-kahani, MR-C303-kissa, MR-C303-kathan-karne, MR-R303-review-296-303

\chapter{Once Upon a Time}
\label{ch:mr-once-upon-a-time}
\hlchaptermodality{303}

\begin{chapteropening}
\textbf{By the end of this chapter:} \emph{I can start a story and say who told it.}

Complete the last lesson of chapter 303: I can start a story and say who told it.
\end{chapteropening}

\section[\texorpdfstring{\mr{एकेकाळी} (ekekāḷī)}{एकेकाळी (ekekāḷī)}]{\mr{एकेकाळी} (ekekāḷī) — once upon a time; at one time}

\label{lesson:MR-C303-ekekali}



\begin{quote}
Before the new one: say the Marathi for at that very moment, then the Marathi for immediately.
\end{quote}

\subsection*{You'll want to know: \mr{एकेकाळी}}
\textbf{\mr{एकेकाळी}} — \emph{ekekāḷī} — \textquotedblleft{}once upon a time; at one time\textquotedblright{}.

\begin{quote}
\mr{एकेकाळी} \mr{या} \mr{डोंगरावर} \mr{एक} \mr{मोठा} \mr{किल्ला} \mr{होता}.
\end{quote}

\begin{quote}
\emph{ekekāḷī yā ḍoṅgrāvar ek moṭhā killā hotā.}
\end{quote}

\begin{quote}
Once upon a time there was a big fort on this hill.
\end{quote}

Built from \mr{एक}, one, and \mr{काळ}, time, with -\mr{ई}: at one time. Stories open with it.

\begin{cousinweb}
\textbf{\mr{एकेकाळी}} (\emph{ekekāḷī}) is built from \emph{ek}, 'one', and \emph{kāḷ}, 'time', which CDIAL lists under Sanskrit \emph{kāla}, 'time' (CDIAL 3084), in the locative: at one time.
\end{cousinweb}

\begin{grammarlens}[title={What you've built}]
One more word for telling what happened.
\end{grammarlens}

\subsection*{Your turn}
\begin{itemize}
\raggedright
  \item \emph{Say it:} \emph{ekekāḷī}
  \item \emph{Say it:} \emph{ekekāḷī}, once more
  \item \emph{Say it:} \emph{ekekāḷī yā ḍoṅgrāvar ek moṭhā killā hotā.} — Once upon a time there was a big fort on this hill.
  \item \emph{Recall:} say the Marathi for at that very moment, then the Marathi for immediately, then say \emph{ekekāḷī} again
\end{itemize}

\begin{tcolorbox}[breakable,colback=teal!4,colframe=teal!35!black,title={Before you move on}]
What does \mr{एकेकाळी} mean? (\textquotedblleft{}Once upon a time; at one time\textquotedblright{}.) Say it once more.
\end{tcolorbox}
\section[\texorpdfstring{\mr{कथा} (kathā)}{कथा (kathā)}]{\mr{कथा} (kathā) — a short story; a story}

\label{lesson:MR-C303-katha}



\begin{quote}
Before the new one: say the Marathi for immediately, then the Marathi for once upon a time.
\end{quote}

\subsection*{You'll want to know: \mr{कथा}}
\textbf{\mr{कथा}} — \emph{kathā} — \textquotedblleft{}a short story; a story\textquotedblright{}.

\begin{quote}
\mr{आजी} \mr{रोज} \mr{रात्री} \mr{मुलांना} \mr{एक} \mr{कथा} \mr{सांगते}.
\end{quote}

\begin{quote}
\emph{ājī roj rātrī mulānnā ek kathā sāṅgte.}
\end{quote}

\begin{quote}
Every night Grandma tells the children a story.
\end{quote}

Feminine.

Plural \mr{कथा}. \mr{गोष्ट} is the everyday word; \mr{कथा} is the written, literary one.

\begin{cousinweb}
\textbf{\mr{कथा}} (\emph{kathā}) is a tatsama, taken from Sanskrit \emph{kathā}, 'a talk, a tale' (Monier-Williams).
\end{cousinweb}

\begin{grammarlens}[title={What you've built}]
One more word for telling what happened.
\end{grammarlens}

\subsection*{Your turn}
\begin{itemize}
\raggedright
  \item \emph{Say it:} \emph{kathā}
  \item \emph{Say it:} \emph{kathā}, once more
  \item \emph{Say it:} \emph{ājī roj rātrī mulānnā ek kathā sāṅgte.} — Every night Grandma tells the children a story.
  \item \emph{Recall:} say the Marathi for immediately, then the Marathi for once upon a time, then say \emph{kathā} again
\end{itemize}

\begin{tcolorbox}[breakable,colback=teal!4,colframe=teal!35!black,title={Before you move on}]
What does \mr{कथा} mean? (\textquotedblleft{}A short story; a story\textquotedblright{}.) Say it once more.
\end{tcolorbox}
\section[\texorpdfstring{\mr{कहाणी} (kahāṇī)}{कहाणी (kahāṇī)}]{\mr{कहाणी} (kahāṇī) — a tale}

\label{lesson:MR-C303-kahani}



\begin{quote}
Before the new one: say the Marathi for once upon a time, then the Marathi for a short story.
\end{quote}

\subsection*{You'll want to know: \mr{कहाणी}}
\textbf{\mr{कहाणी}} — \emph{kahāṇī} — \textquotedblleft{}a tale\textquotedblright{}.

\begin{quote}
\mr{त्या} \mr{जुन्या} \mr{घराची} \mr{कहाणी} \mr{फार} \mr{दुःखाची} \mr{आहे}.
\end{quote}

\begin{quote}
\emph{tyā junyā gharācī kahāṇī phār duḥkhācī āhe.}
\end{quote}

\begin{quote}
The tale of that old house is very sad.
\end{quote}

Feminine. Plural: \mr{कहाण्या}, \emph{kahāṇyā}.

Plural \mr{कहाण्या}.

\begin{cousinweb}
\textbf{\mr{कहाणी}} (\emph{kahāṇī}) is a tadbhava: CDIAL lists Marathi \emph{kahāṇī}, 'a tale', under Sanskrit \emph{kathānaka}, 'a little tale' (CDIAL 2705).
\end{cousinweb}

\begin{grammarlens}[title={What you've built}]
One more word for telling what happened.
\end{grammarlens}

\subsection*{Your turn}
\begin{itemize}
\raggedright
  \item \emph{Say it:} \emph{kahāṇī}
  \item \emph{Say it:} \emph{kahāṇī}, once more
  \item \emph{Say it:} \emph{tyā junyā gharācī kahāṇī phār duḥkhācī āhe.} — The tale of that old house is very sad.
  \item \emph{Recall:} say the Marathi for once upon a time, then the Marathi for a short story, then say \emph{kahāṇī} again
\end{itemize}

\begin{tcolorbox}[breakable,colback=teal!4,colframe=teal!35!black,title={Before you move on}]
What does \mr{कहाणी} mean? (\textquotedblleft{}A tale\textquotedblright{}.) Say it once more.
\end{tcolorbox}
\section[\texorpdfstring{\mr{किस्सा} (kissā)}{किस्सा (kissā)}]{\mr{किस्सा} (kissā) — an anecdote}

\label{lesson:MR-C303-kissa}



\begin{quote}
Before the new one: say the Marathi for a short story, then the Marathi for a tale.
\end{quote}

\subsection*{You'll want to know: \mr{किस्सा}}
\textbf{\mr{किस्सा}} — \emph{kissā} — \textquotedblleft{}an anecdote\textquotedblright{}.

\begin{quote}
\mr{बाबांनी} \mr{त्यांच्या} \mr{लहानपणीचा} \mr{एक} \mr{मजेदार} \mr{किस्सा} \mr{सांगितला}.
\end{quote}

\begin{quote}
\emph{bābānnī tyāñcyā lahānapṇīcā ek majedār kissā sāṅgitlā.}
\end{quote}

\begin{quote}
Dad told a funny anecdote from his childhood.
\end{quote}

Masculine. Plural: \mr{किस्से}, \emph{kisse}.

Plural \mr{किस्से}.

\begin{cousinweb}
\textbf{\mr{किस्सा}} (\emph{kissā}) is a loan from Persian \emph{qissa}, usually traced to Arabic \emph{qiṣṣa}, 'a story, a narrative'.
\end{cousinweb}

\begin{grammarlens}[title={What you've built}]
One more word for telling what happened.
\end{grammarlens}

\subsection*{Your turn}
\begin{itemize}
\raggedright
  \item \emph{Say it:} \emph{kissā}
  \item \emph{Say it:} \emph{kissā}, once more
  \item \emph{Say it:} \emph{bābānnī tyāñcyā lahānapṇīcā ek majedār kissā sāṅgitlā.} — Dad told a funny anecdote from his childhood.
  \item \emph{Recall:} say the Marathi for a short story, then the Marathi for a tale, then say \emph{kissā} again
\end{itemize}

\begin{tcolorbox}[breakable,colback=teal!4,colframe=teal!35!black,title={Before you move on}]
What does \mr{किस्सा} mean? (\textquotedblleft{}An anecdote\textquotedblright{}.) Say it once more.
\end{tcolorbox}
\section[\texorpdfstring{\mr{कथन} \mr{करणे} (kathan karṇe)}{कथन करणे (kathan karṇe)}]{\mr{कथन} \mr{करणे} (kathan karṇe) — to narrate; to recount}

\label{lesson:MR-C303-kathan-karne}



\begin{quote}
Before the new one: say the Marathi for a tale, then the Marathi for an anecdote.
\end{quote}

\subsection*{You'll want to know: \mr{कथन} \mr{करणे}}
\textbf{\mr{कथन} \mr{करणे}} — \emph{kathan karṇe} — \textquotedblleft{}to narrate; to recount\textquotedblright{}.

\begin{quote}
\mr{साक्षीदाराने} \mr{घडलेली} \mr{घटना} \mr{न्यायालयात} \mr{कथन} \mr{केली}.
\end{quote}

\begin{quote}
\emph{sākṣīdārāne ghaḍlelī ghaṭnā nyāyālayāt kathan kelī.}
\end{quote}

\begin{quote}
The witness recounted the event in court.
\end{quote}

Past, with -\mr{ने} (he …): \mr{त्याने} \mr{कथन} \mr{केलं}, \emph{tyāne kathan kelaṁ}.

Formal; in speech \mr{सांगणे}. What is told agrees: \mr{घटना} \mr{कथन} \mr{केली}.

\begin{cousinweb}
\textbf{\mr{कथन} \mr{करणे}} (\emph{kathan karṇe}) joins \emph{kathan}, 'a telling', a tatsama taken from Sanskrit \emph{kathana}, 'narration' (Monier-Williams), to \emph{karṇe}, 'to do'.
\end{cousinweb}

\begin{grammarlens}[title={What you've built}]
One more word for telling what happened.
\end{grammarlens}

\subsection*{Your turn}
\begin{itemize}
\raggedright
  \item \emph{Say it:} \emph{kathan karṇe}
  \item \emph{Say it:} \emph{kathan karṇe}, once more
  \item \emph{Say it:} \emph{sākṣīdārāne ghaḍlelī ghaṭnā nyāyālayāt kathan kelī.} — The witness recounted the event in court.
  \item \emph{Recall:} say the Marathi for a tale, then the Marathi for an anecdote, then say \emph{kathan karṇe} again
\end{itemize}

\begin{tcolorbox}[breakable,colback=teal!4,colframe=teal!35!black,title={Before you move on}]
What does \mr{कथन} \mr{करणे} mean? (\textquotedblleft{}To narrate; to recount\textquotedblright{}.) Say it once more.
\end{tcolorbox}
\section[(dialogue)]{Looking back: from A Crash on the Road to Once Upon a Time}

\label{lesson:MR-R303-review-296-303}



\begin{quote}
Say the Marathi for to collide and to slam into.
\end{quote}

\subsection*{Your turn}
Say these aloud:

\begin{itemize}
\raggedright
  \item to collide, to slam into, unconscious, to recover, to be missed
  \item to barge in, to break open, to snatch, to loot, to fleece
  \item to scream, commotion, havoc, a mob, to whisper
  \item to search a place, a scene report, to confiscate, to confess, to file
  \item to threaten, a bribe, a punishment, to flee, to shut in
  \item to be seen, to be found, to come across, to peep, to observe closely
  \item to be startled, to stop short, a moment, at that very moment, immediately
  \item once upon a time, a short story, a tale, an anecdote, to narrate
\end{itemize}

\begin{tcolorbox}[breakable,colback=teal!4,colframe=teal!35!black,title={Before you move on}]
To collide? (\textbf{\mr{धडकणे}}, \emph{dhaḍakṇe}.) To slam into? (\textbf{\mr{आदळणे}}, \emph{ādaḷṇe}.) Unconscious? (\textbf{\mr{बेशुद्ध}}, \emph{beśuddha}.) To recover? (\textbf{\mr{सावरणे}}, \emph{sāvarṇe}.) To be missed? (\textbf{\mr{हुकणे}}, \emph{hukṇe}.) To barge in? (\textbf{\mr{घुसणे}}, \emph{ghusṇe}.) To break open? (\textbf{\mr{फोडणे}}, \emph{phoḍṇe}.) To snatch? (\textbf{\mr{हिसकावणे}}, \emph{hiskāvṇe}.) To loot? (\textbf{\mr{लुटणे}}, \emph{luṭṇe}.) To fleece? (\textbf{\mr{लुबाडणे}}, \emph{lubāḍṇe}.) To scream? (\textbf{\mr{किंचाळणे}}, \emph{kiñcāḷṇe}.) Commotion? (\textbf{\mr{गोंधळ}}, \emph{gondhaḷ}.) Havoc? (\textbf{\mr{हाहाकार}}, \emph{hāhākār}.) A mob? (\textbf{\mr{जमाव}}, \emph{jamāv}.) To whisper? (\textbf{\mr{कुजबुजणे}}, \emph{kujbujṇe}.) To search a place? (\textbf{\mr{झडती} \mr{घेणे}}, \emph{jhaḍtī gheṇe}.) A scene report? (\textbf{\mr{पंचनामा}}, \emph{pañcnāmā}.) To confiscate? (\textbf{\mr{जप्त} \mr{करणे}}, \emph{japta karṇe}.) To confess? (\textbf{\mr{कबूल} \mr{करणे}}, \emph{kabūl karṇe}.) To file? (\textbf{\mr{दाखल} \mr{करणे}}, \emph{dākhal karṇe}.) To threaten? (\textbf{\mr{धमकावणे}}, \emph{dhamkāvṇe}.) A bribe? (\textbf{\mr{लाच}}, \emph{lāc}.) A punishment? (\textbf{\mr{शिक्षा}}, \emph{śikṣā}.) To flee? (\textbf{\mr{पळ} \mr{काढणे}}, \emph{paḷ kāḍhṇe}.) To shut in? (\textbf{\mr{कोंडणे}}, \emph{koṇḍṇe}.) To be seen? (\textbf{\mr{दिसणे}}, \emph{disṇe}.) To be found? (\textbf{\mr{सापडणे}}, \emph{sāpaḍṇe}.) To come across? (\textbf{\mr{आढळणे}}, \emph{āḍhaḷṇe}.) To peep? (\textbf{\mr{डोकावणे}}, \emph{ḍokāvṇe}.) To observe closely? (\textbf{\mr{निरखणे}}, \emph{nirakhṇe}.) To be startled? (\textbf{\mr{दचकणे}}, \emph{dacakṇe}.) To stop short? (\textbf{\mr{थबकणे}}, \emph{thabakṇe}.) A moment? (\textbf{\mr{क्षण}}, \emph{kṣaṇ}.) At that very moment? (\textbf{\mr{तेवढ्यात}}, \emph{tevaḍhyāt}.) Immediately? (\textbf{\mr{ताबडतोब}}, \emph{tābaḍtob}.) Once upon a time? (\textbf{\mr{एकेकाळी}}, \emph{ekekāḷī}.) A short story? (\textbf{\mr{कथा}}, \emph{kathā}.) A tale? (\textbf{\mr{कहाणी}}, \emph{kahāṇī}.) An anecdote? (\textbf{\mr{किस्सा}}, \emph{kissā}.) To narrate? (\textbf{\mr{कथन} \mr{करणे}}, \emph{kathan karṇe}.)
\end{tcolorbox}
